\documentclass[11pt]{article}
\usepackage{amsmath,amssymb}
\title{The latest clearance report with a safe rolling approach}
\begin{document}
\maketitle

Waiting at the stop line is safe but need not be the fastest way to
finish after clearance arrives. A vehicle can arrange to arrive at the
decision with positive speed while preserving a full stop if clearance
is withheld. The best speed can be strictly below the fastest safe
approach speed: a faster vehicle needs to be farther from the line
in order to retain its stop.

\paragraph{Interface.}
A civil corridor vehicle starts at rest at zero, moves forward only,
and has an acceleration magnitude bound \(a>0\), with no binding
speed cap. Its stop line is \(\ell>0\); the target is \(L>\ell\)
and requires terminal rest. A truthful report arrives at known time
\(q\ge0\), distinguishing a crossing that is physically clear and
cannot be reoccupied from a crossing that must not be entered.
Before that report, controls must be identical. In the latter
case (or with no report) the vehicle must be able to stop no later
than \(\ell\), without crossing. The line already includes any
space required for clearance. No risk is inferred from the
software assignment itself; the source of the physical report is
an explicit premise.

\paragraph{Exact decision at the cut.}
For fixed speed \(v\) at \(q\), the maximum distance traveled from
rest is the integral of
\(\min(at,v+a(q-t))\), namely
\(a q^2/4+qv/2-v^2/(4a)\). The minimum is \(v^2/(2a)\),
by waiting and using the last part of the interval to accelerate.
The two profiles interpolate to any intermediate distance. To
retain a stop at the line, position \(x\) must satisfy
\(x+v^2/(2a)\le\ell\). The stopping coordinate is nondecreasing
along every admissible forward prefix:
its derivative is \(v(1+\dot v/a)\ge0\) almost everywhere.
Consequently satisfying the bound at \(q\) ensures it beforehand.
The exact speed range and best position for each speed are
\[
  0\le v\le\min(aq,\sqrt{a\ell}), \qquad
  x_*(v)=\min\left(\frac{a q^2}{4}+\frac{qv}{2}
                     -\frac{v^2}{4a},\ \ell-\frac{v^2}{2a}\right).
\]
This is attainable with a single prefix for both branches.
Without clearance, brake to a stop. With clearance, the minimum
additional time at this prefix is
\[
   \frac{2\sqrt{a(L-x_*(v))+v^2/2}-v}{a}.
\]
Indeed the maximum displacement starting at \(v\) and ending at
rest in time \(t\ge v/a\) is the integral of
\(\min(v+au,a(t-u))\), equal to
\(at^2/4+vt/2-v^2/(4a)\). Solving at distance \(L-x_*(v)\) gives
the expression and the attaining velocity profile. Taking the
minimum over the exact speed interval solves the original problem.
When \(q=0\) the range collapses to zero.

\paragraph{A closed-form case.}
If \(q\ge 2\sqrt{\ell/a}\), every boundary state
\((\ell-v^2/(2a),v)\) in the speed range is attainable. For example,
time the fastest rest-to-rest approach to \(\ell\) so that its
deceleration passes through that state at \(q\). On the boundary the
objective becomes
\((2\sqrt{a(L-\ell)+v^2}-v)/a\).
It is minimized at \(\min(\sqrt{a(L-\ell)/3},\sqrt{a\ell})\).
For \(L\le4\ell\) it equals \(\sqrt{3(L-\ell)/a}\), compared
with \(2\sqrt{(L-\ell)/a}\) at rest on the line. For \(L>4\ell\)
it equals \(2\sqrt{L/a}-\sqrt{\ell/a}\). The calculation of the
boundary also applies at an earlier \(q\) whenever its minimizing
state can already be reached; the one-dimensional optimization
covers the remaining early case.

\paragraph{Replay and comparison.}
Use \(a=\ell=1\), \(L=28/25\) and \(q=11/5\). The exact prefix
can wait \(2/5\), accelerate for one, and brake for \(4/5\).
At the cut it is at \(49/50\) with speed \(1/5\); it can stop
exactly at the line if the report is absent. On clearance,
accelerating for \(1/5\) and then braking for \(2/5\) reaches the
target at \(14/5\). Waiting stopped at the line finishes only
at about \(2.89282\). The replay checks the exact piecewise paths,
an independent sampled velocity LP at the named speed, and
81 numerical comparisons of the reduced scalar objective.

The example uses the earlier inherited-state motion calculation
and the native assignment replay's distinction between a grant
and physical exit. It does not supply a real observation channel.
An equally informed direct trajectory optimizer can choose the
same prefix. No claim is made for binding speed limits, delayed
or incorrect reports, shared risk estimates, or a deployed
crossing controller. This is one scoped model result, not I1 acceptance.
\end{document}
